\section{Discussion}
\label{sec:discussion}

\subsection{FT-Transformer as a Financial Encoder}
\label{sec:disc_encoder}

Against default-hyperparameter XGBoost, the tuned FT-Transformer is directionally positive on all four distress outcomes ($+0.015$ to $+0.075$ AUROC), with statistically significant gains on \texttt{stock\_decline} and \texttt{earnings\_restate}. When the comparison is made like-for-like---both the baselines and the FT-Transformer Optuna-tuned on the same expanding-window protocol---the FT-Transformer meets the gate (3/4 wins): \texttt{earnings\_restate} ($+0.037$), \texttt{bankruptcy} ($+0.041$), and \texttt{sec\_enforcement} ($+0.016$), with a sub-threshold lead on \texttt{stock\_decline} ($+0.009$). The advantage is uneven, however: simpler models lead on \texttt{stock\_decline} (LR-full) and \texttt{sec\_enforcement} (LR-traditional), consistent with findings from \citet{gorishniy2021revisiting} that the relative advantage of attention-based models over tree ensembles is dataset-dependent and sensitive to tuning.

The calibration story provides additional nuance. On \texttt{stock\_decline}---the highest-prevalence and most statistically reliable outcome---the FT-Transformer achieves substantially better calibration than both default and tuned XGBoost (ECE 0.038 vs.\ 0.152 for tuned XGBoost). On low-base-rate outcomes, XGBoost achieves near-zero ECE by predicting near-zero probabilities for virtually all samples, which is technically well-calibrated but uninformative for risk management. The FT-Transformer produces a wider spread of predicted probabilities that enables practical threshold-setting and risk differentiation. For the downstream JEPA use case, where the encoder must produce well-structured continuous representations rather than binary classifications, this distributional richness matters more than marginal AUROC.

\subsection{Self-Supervised Pretraining}
\label{sec:disc_ssl}

SSL pretraining via masked feature reconstruction provides meaningful gains, with the most dramatic improvement on \texttt{sec\_enforcement} ($+0.093$ AUROC at $r = 0.50$). This outcome has the lowest base rate and the weakest supervised signal---precisely the regime where pretraining should help most, since the encoder can learn financial structure from the full unlabeled dataset rather than relying on scarce positive examples.

The absence of a single optimal mask ratio across outcomes suggests that different outcomes depend on different structural properties of the financial data. Low mask ratios (0.15) may encourage the encoder to learn fine-grained local dependencies between related features (e.g., within the balance sheet), while high mask ratios (0.50) force it to learn more global, abstract relationships. The latter appears more useful for outcomes like \texttt{sec\_enforcement} that may depend on higher-order patterns not captured by any individual financial ratio.

\subsection{Implications for Study~1}
\label{sec:disc_study1}

The FT-Transformer meets the AUROC gate against fully Optuna-tuned baselines (3/4 wins), and walk-forward validation confirms the advantage is temporally stable on \texttt{earnings\_restate}. We adopt the FT-Transformer as the financial modality encoder for Study~1. Beyond the gate itself, three properties reinforce this choice:

\begin{enumerate}[leftmargin=*]
  \item \textbf{Representation quality over marginal discrimination.} Study~1 uses the encoder's latent representations for temporal prediction, not its classification outputs. SSL pretraining improves downstream prediction even though the pretraining objective (feature reconstruction) has no direct relationship to distress labels---demonstrating that the FT-Transformer's latent space captures meaningful financial structure. If the latent space is rich enough to support feature reconstruction, it may also support temporal prediction (Study~1) and cross-modal alignment (Study~2). XGBoost, as a non-differentiable ensemble, cannot serve this role.

  \item \textbf{Calibration and distributional properties.} On \texttt{stock\_decline}, the FT-Transformer produces substantially better-calibrated probabilities (ECE 0.038 vs.\ 0.152 for tuned XGBoost). On low-base-rate outcomes, the FT-Transformer's wider probability spread enables practical risk stratification, whereas XGBoost's near-degenerate predictions offer no discriminatory value. For a downstream architecture that consumes continuous representations, well-structured output distributions are essential.

  \item \textbf{Robustness of the headline advantage.} Walk-forward validation across seven expanding windows shows the FT-Transformer's edge on \texttt{earnings\_restate}---the highest-signal of the rare outcomes---holds in every fold, including the two COVID-era windows, rather than being an artifact of the single 2020--2023 test split. The bankruptcy advantage is regime-sensitive (it reverses in the COVID folds), which we treat as a caveat rather than a basis for the decision.
\end{enumerate}

Additionally, SSL pretraining remains clearly beneficial: no single mask ratio dominates, but $r = 0.50$ yields the strongest gains on the most challenging outcomes (\texttt{sec\_enforcement}, \texttt{earnings\_restate}). The grid search (\Cref{fig:sweep_heatmap}) shows moderate hyperparameter sensitivity (up to 0.05 mean AUROC), with learning rate the most influential factor.

\subsection{Limitations}
\label{sec:limitations}

Several limitations should be noted. Although the FT-Transformer clears the AUROC gate against fully tuned baselines (3/4 wins), the margins are uneven and two of the three wins warrant caution. The \texttt{sec\_enforcement} result is particularly unstable: AUROC varies substantially across runs and seeds (multi-seed std $= 0.012$), a traditional-ratio logistic regression outperforms every other model on it, and the FT-Transformer wins only 3 of 7 walk-forward folds---all reflecting the extremely low base rate (AUPRC $< 0.02$). The \texttt{bankruptcy} outcome shows similar noise (multi-seed std $= 0.013$; bootstrap CI includes zero) and is regime-sensitive, reversing in both COVID-era walk-forward folds. While multi-seed evaluation (3 seeds) confirms modest variance for the two higher-prevalence outcomes (std $< 0.01$), the rare-event outcomes remain noisy; the most robust advantage is on \texttt{earnings\_restate}, which holds across all seven walk-forward folds. The walk-forward numbers themselves use a normalization scaler fit on data up to 2017 and shared across folds, so the two earliest folds incur mild feature-level look-ahead; refitting the scaler within each fold is the strictly leak-free protocol and is the one correction we would apply before treating the early-fold magnitudes as exact (the fold-level FT-vs-XGBoost sign, which both models share, is unaffected). Finally, the test period (2020--2023) spans the COVID-19 market disruption, which provides a valuable out-of-distribution stress test but also means the test distribution may differ systematically from the training distribution.

\subsection{Related Work}
\label{sec:related}

\paragraph{Deep learning on tabular data.} Whether deep learning can match tree-based ensembles on tabular tasks remains contested. \citet{gorishniy2021revisiting} showed that the FT-Transformer matches or exceeds XGBoost on a range of benchmarks, while \citet{grinsztajn2022tree} and \citet{shwartzziv2022tabular} found that trees remain competitive or superior on many datasets. Other attention-based approaches include TabNet \citep{arik2021tabnet}, which integrates feature selection with attentive learning. Our work contributes a domain-specific validation on financial XBRL data, where extreme tails, systematic missingness, and temporal structure distinguish the setting from generic tabular benchmarks.

\paragraph{Financial distress prediction.} Machine learning for bankruptcy and distress prediction has a long history, from \citet{altman1968financial}'s Z-score to modern ensemble methods. \citet{barboza2017machine} benchmarked random forests and boosting against traditional statistics, establishing ML superiority. \citet{mai2019deep} applied deep learning to textual disclosures alongside accounting ratios. More recently, \citet{korangi2023transformer} applied Transformers to structured financial panel data for corporate default prediction, and \citet{chen2022predicting} demonstrated the value of high-dimensional XBRL-tagged features over traditional small variable sets. Our work differs in using raw XBRL features extracted directly from SEC EDGAR (rather than pre-curated Compustat fields) and in validating self-supervised pretraining in this domain.

\paragraph{Self-supervised pretraining for tabular data.} Adapting SSL to tabular data is a recent research direction. VIME \citep{yoon2020vime} introduced mask estimation and value imputation as pretext tasks for tabular features. SCARF \citep{bahri2022scarf} applied contrastive learning via random feature corruption, while SubTab \citep{ucar2021subtab} used feature subsetting to create multiple views. Our masked feature reconstruction objective is closest to VIME but uses a simpler masking strategy (zero-replacement rather than learned corruption) and reconstructs from a global CLS representation rather than per-feature outputs.
